\scanpage{25}

les $Dt_i=0$. \deleted{\unsure{Alors}} Prenons une
\uline{base} de $\Omega_{K/k}$ sur $K$, \unsure{soit} $(\omega_i)_{i\in I}$,
donc $\Omega_{K/k}\otimes_K A\simeq A^{(I)}\subset A^I$. En tant
que dérivation dans $A^I$, $D$ est donnée
par ses composantes $D_i$, qui sont
\unsure{uniq}$^t$ déterminées
\[
D_i\Bigl(\sum_\lambda a_\lambda t^\lambda\Bigr)
=\sum_\lambda \unsure{d^i_{K/k}|d_{K/k}^{\,i}}(a_\lambda)t^\lambda
\]
Il reste à prouver qu'on peut trouver
des $f=\sum a_\lambda t^\lambda$ telles que les
\[
\bigl(D_i(\unsure{f|\varphi})\bigr)_{i\in I}\notin A^{(I)}.
\]
\deleted{Il suffit}
\inserted{Choisissons}
\deleted{pour ceci \unsure{que} \inserted{\unsure{pour avoir}} \unsure{les}
$D_i(a_\lambda)$ \unsure{pour} $\unsure{\forall}$}
\deleted{de \unsure{trouver}}
une \uline{suite infinie} $i_0,i_1,\ldots,i_n$
d'indices distincts de $I$, \deleted{telle que} pour $\unsure{\forall}$
\deleted{$D_{i_n}(a_n)\ne0$}
\unsure{n} il existe $a_n\in K$ tel que
\deleted{$\unsure{D_{i_n}(a_n)\ne0}$},
\deleted{\unsure{(ceci est possible}}
$d_{K/k}a_n^{\unsure{i_n}}\ne0$, il suffit alors de
prendre $f=\sum a_nt^n$. .\,.\,.\,.\,.

\mspage{24}
